\chapter{Some definitions concerning curves, surfaces, and regions}
\label{appendix: Some definitions concerning curves, surfaces, and regions}

\begin{enumerate}

\item
A curve $C\colon \vec{r} = \vec{r}(t)$, $t\in[a,b]$ is a \emph{smooth curve} if $d\vec{r}(t)/dt$ exists, is continuous and is non-zero on $[a,b]$. 

\item
A curve $C$ is \emph{piecewise smooth} if it is continuous and consists of a finite number of smooth curves.

\item
A curve $C$ is \emph{simple} if it is not self-intersecting.

\item
A curve $C$ is \emph{regular} if it is simple and piecewise smooth.

\item
A surface $S$ is \emph{smooth} if at each point of $S$ there exists a unique unit normal vector which varies continuously along $S$. $S$ is \emph{piecewise smooth} if it consists of a finite number of smooth pieces. 

\item
A smooth surface is \emph{orientable} if it is two-side, i.e. it is possible on the entire surface to identify a positive normal direction in a unique and continuous way. 

\item
A surface is \emph{regular} if it is orientable and piecewise smooth. 

\item
A set of points $S$ is any collection of points in $E_3$. The \emph{complementary of $S$} is the set of all points in $E_3$ that are not in $S$. 

\item
$S$ is \emph{bounded} if all points of $S$ lie within a sphere of finite radius. Otherwise, $S$ is \emph{unbounded}. 

\item
A $\delta$ neighborhood of a point $P$ is the set of points lying inside a sphere of radius $\delta$ centered at $P$.

\item
A point $P\in S$ is an \emph{interior point} of $S$ if there is a neighborhood of $P$ lying entirely in $S$. 

\item
A set $S$ is \emph{open} if every point in $S$ is an interior point. 

\item
A set $S$ is \emph{connected} if any two points in $S$ can be connected by a finite-length path lying entirely in $S$. 

\item
A \emph{domain} (or open region) is an open connected set. 

\item
A point $P$ (not necessary in $S$) is a \emph{limit point (or accumulation pint) if every neighborhood of $P$} (no matter how small it is) contains points of $S$ distinct from $P$. 

\item
A set is \emph{closed} if it contains all of its limit points. 

\item
A set is \emph{compact} if it is both closed and bounded. 

\item
The \emph{closure} of $S$, denoted by $\bar{S}$, is the union of $S$ and all of its limit points. 

\item
A \emph{boundary point} of $S$ is a limit point which is not an interior point. The boundary of $S$ is denoted by $\partial S$. The interior of $S$ is denoted by $\mathring{S}$. 

\item
A region is a domain together with all, some , or none of its boundary points. In other words, a region may be neither open nor closed. 

\item
A region $R$ is \emph{simply-connected} if every closed curve \emph{lying entirely in $R$} can be continuously deformed to a point without leaving $R$. Otherwise, it is \emph{multiply-connected}. 


\end{enumerate}